
\chapter{Quadrupole and $L_{GW}$}\label{ap:L_GW}

In this appendix we will go through the details of finding the mass quadrupole components, explain the approximations assumed and finding the $L_{GW}$.
Lets start with the definitions, as we saw from Fig. \ref{fig:binary} the position are:

\begin{align}
    \textbf{r}_1 &= -\frac{\mu a}{m_1}(cos(\theta),sin(\theta),0),\\
    \textbf{r}_2 &= \frac{\mu a}{m_2}(cos(\theta),sin(\theta),0).\\
\end{align}

\noindent We used $||\textbf{r}_1||= m_2 a/(m_1+m_2)=\mu a/m_1$ with $\mu=m_1 m_2 /(m_1+m_2)$.
\noindent The first quadrupole component will be:
\begin{align}
    Q_{xx}&=m_1r_{x1}r_{x1}+m_2r_{x2}r_{x2}=m_1\frac{\mu^2 a^2}{m_1^2}cos^2(\theta)+m_2\frac{\mu^2 a^2}{m_2^2}cos^2(\theta),\\
    Q_{xx}&=\mu^2 a^2cos^2(\theta)(\frac{1}{m_1}+\frac{1}{m_2}),\\
    Q_{xx}&=\mu a^2cos^2(\theta),\\
    Q_{xx}&=\frac{\mu a^2}{2}(cos(2\theta)+1).\\
\end{align}
\noindent Similarly for $Q_{yy}$:
\begin{align}
    Q_{yy}&=m_1r_{y1}r_{y1}+m_2r_{y2}r_{y2}=m_1\frac{\mu^2 a^2}{m_1^2}sin^2(\theta)+m_2\frac{\mu^2 a^2}{m_2^2}sin^2(\theta),\\
    Q_{yy}&=\mu^2 a^2sin^2(\theta)(\frac{1}{m_1}+\frac{1}{m_2}),\\
    Q_{yy}&=\mu a^2sin^2(\theta),\\
    Q_{yy}&=\frac{\mu a^2}{2}(-cos(2\theta)+1).\\
\end{align}
\noindent Finally for $Q_{xy}$
\begin{align}
    Q_{xy}&=m_1r_{x1}r_{y1}+m_2r_{x2}r_{y2}=m_1\frac{\mu^2 a^2}{m_1^2}sin(\theta)cos(\theta)+m_2\frac{\mu^2 a^2}{m_2^2}sin(\theta)cos(\theta),\\
    Q_{xy}&=\mu^2 a^2sin(\theta)cos(\theta)(\frac{1}{m_1}+\frac{1}{m_2}),\\
    Q_{xy}&=\mu a^2sin(\theta)cos(\theta),\\
    Q_{xy}&=\frac{\mu a^2}{2}sin(2\theta).\\
\end{align}
\noindent Now we need the third derivate with respect to time. 
We will assume two things, $\ddot\theta<<1$ and $\dot a<<1$.
Effectively we are saying that the change in angular velocity and the change in the radius is slow which is referred as an adiabatic process. 
To see how good the approximation is we can compare the orbit time scale with the inspiral timescale.
Which in fact is the situation, therefore is a good approximation until the merger itself.
Therefore it is easy to see that:
\begin{align}
    \dddot Q_{xx}&=4\mu a^2\dot\theta^3sin(2\theta)\\
    \dddot Q_{yy}&=-4\mu a^2\dot\theta^3sin(2\theta)\\
    \dddot Q_{xy}&=-4\mu a^2\dot\theta^3cos(2\theta)\\
\end{align}
Now we meed to calculate the following term remembering that we are using einstein notation for the repeated indices and that $Q_{xy}=Q_{yx}$:
\begin{align}
    \dddot Q_{ij}\dddot Q_{ij} - \dddot Q_{ii}\dddot Q_{jj}&=\dddot Q_{xx}\dddot Q_{xx}+\dddot Q_{xy}\dddot Q_{xy}+\dddot Q_{yx}\dddot Q_{yx}+\dddot Q_{yy}\dddot Q_{yy}\\
    &-\dddot Q_{xx}\dddot Q_{xx}-\dddot Q_{yy}\dddot Q_{xx}-\dddot Q_{xx}\dddot Q_{yy}-\dddot Q_{yy}\dddot Q_{yy}\\
    &=\dddot Q_{xx}^2+2\dddot Q_{xy}^2+\dddot Q_{yy}^2-\dddot Q_{xx}^2-2\dddot Q_{yy}\dddot Q_{xx}-\dddot Q_{yy}^2\\
    &=\dddot Q_{xx}^2+2\dddot Q_{xy}^2+\dddot Q_{yy}^2-(\dddot Q_{xx}+\dddot Q_{yy})^2\\
    &=16\mu^2 a^4 \dot\theta^6(sin^2(2\theta)+2cos^2(2\theta)+sin^2(2\theta))\\
    &=32\mu^2 a^4 \dot\theta^6=\frac{32G^3m_1^2m_2^2M}{a^5}
\end{align}
Finally:
\begin{equation}
    L_{GW}=\frac{G}{5c^5}(\dddot Q_{ij}\dddot Q_{ij} - \dddot Q_{ii}\dddot Q_{jj})=\frac{32G^4 m_1^2 m_2^2M}{c^5a^5}
\end{equation}
